\chapter{Homotopia, nakrycia i komórkowe modele przestrzeni}
\label{ch:topologia-algebraiczna-podstawy}
\index{grupa podstawowa}

W Rozdziale~\ref{ch:algebra-homologiczna} kompleks łańcuchowy
był ciągiem modułów i homomorfizmów. Aby zbudować taki
kompleks z przestrzeni, nauczymy się rozkładać przestrzeń
na proste części. Równolegle zapytamy, kiedy dwa odwzorowania
można ciągle deformować jedno w drugie. Są to pytania
topologiczne: gładkość nie jest tu potrzebna.

Zaczniemy od dróg i nakryć, następnie od sympleksów
przejdziemy do CW-kompleksów. Homologia tych obiektów
pojawi się w następnym rozdziale. Homotopii odwzorowań
nie należy mylić z homotopią \emph{łańcuchową}
z Definicji~\ref{def:homotopia-lancuchowa};
ta druga jest algebraicznym odpowiednikiem pierwszej.

\section{Homotopie, drogi i grupa podstawowa}
\label{sec:homotopie-pi1}

W całym rozdziale $I=[0,1]$. Parametr $t\in I$
można interpretować jako czas ciągłej deformacji.
Obrazy pośrednie mogą być bardzo różne, lecz
ciągłość wyklucza skoki i rozdarcia.

\begin{definicja}[Homotopia i równoważność homotopijna]
\index{homotopia i rownowaznosc homotopijna@Homotopia i równoważność homotopijna}
\label{def:homotopia-topologiczna}
Dwa ciągłe odwzorowania $f,g:X\to Y$ są
\emph{homotopijne}, zapisujemy $f\simeq g$,
jeśli istnieje ciągłe $H:X\times I\to Y$ z
$H(x,0)=f(x)$ i $H(x,1)=g(x)$.
Jeśli $f|_A=g|_A$ oraz $H(a,t)=f(a)$
dla $a\in A\subseteq X$ i $t\in I$,
mówimy o homotopii \emph{względem $A$}.
Przestrzenie $X,Y$ są \emph{homotopijnie
równoważne}, gdy istnieją ciągłe $f:X\to Y$,
$g:Y\to X$ z $gf\simeq\id_X$ i
$fg\simeq\id_Y$.
\end{definicja}

Homeomorfizm ma prawdziwą odwrotność.
Równoważność homotopijna dopuszcza odwrotność
z dokładnością do deformacji. Dysk i punkt
nie są homeomorficzne, ale są homotopijnie
równoważne.

\begin{definicja}[Retrakt deformacyjny]
\index{retrakt deformacyjny@Retrakt deformacyjny}
\label{def:retrakt-deformacyjny}
Podprzestrzeń $A\subseteq X$ jest
\emph{silnym retraktem deformacyjnym} $X$,
jeśli istnieje ciągłe $H:X\times I\to X$ takie, że
\[
 H(x,0)=x,\qquad H(x,1)\in A,\qquad
 H(a,t)=a\quad(a\in A).
\]
Odwzorowanie $r(x):=H(x,1)$ jest wówczas
retrakcją $X\to A$. Przestrzeń homotopijnie
równoważną punktowi nazywamy \emph{ściągalną}.
\end{definicja}

\begin{przyklad}[Promieniowa deformacja]
\label{prz:retrakcja-pierscien}
Dla $x\in\R^n\setminus\{0\}$ wzór
\[
 H(x,t)=\left((1-t)+\frac{t}{\|x\|}\right)x
\]
daje silny retrakt deformacyjny
$\R^n\setminus\{0\}$ na $S^{n-1}$.
Współczynnik jest dodatni, więc tor
nie przechodzi przez usunięty początek.
Dla $\|x\|=1$ punkt pozostaje nieruchomy,
a dla $t=1$ otrzymujemy $x/\|x\|$.
Każdy niepusty zbiór wypukły $C\subseteq\R^N$ jest ściągalny:
dla wybranego $c_0\in C$ stosujemy
$H(x,t)=(1-t)x+tc_0$.
\end{przyklad}

Do badania dziur potrzebujemy dróg,
których końce pozostają nieruchome w czasie deformacji.

\begin{definicja}[Drogi, pętle i homotopia dróg]
\index{drogi petle i homotopia drog@Drogi, pętle i homotopia dróg}
\label{def:droga-petla}
\emph{Droga} w $X$ to ciągłe
$\alpha:I\to X$; jej początek i koniec
to $\alpha(0)$ i $\alpha(1)$.
Gdy oba są równe $x_0$, nazywamy ją
\emph{pętlą opartą w $x_0$}.
Drogi $\alpha,\beta$ o tych samych końcach
są homotopijne \emph{względem końców},
jeśli łączy je homotopia $H:I\times I\to X$
z $H(0,t)=\alpha(0)$ i
$H(1,t)=\alpha(1)$ dla każdego $t$.
Gdy $\alpha(1)=\beta(0)$, określamy
\[
 (\alpha*\beta)(s)=
 \begin{cases}
  \alpha(2s),&0\leq s\leq\tfrac12,\\
  \beta(2s-1),&\tfrac12\leq s\leq1,
 \end{cases}
 \qquad \bar\alpha(s)=\alpha(1-s).
\]
Symbol $\alpha*\beta$ oznacza: najpierw $\alpha$,
potem $\beta$.
\end{definicja}

Homotopia względem końców jest relacją równoważności.
Stała deformacja daje zwrotność, zamiana $t$ na $1-t$ daje
symetrię, a sklejenie dwóch homotopii na połówkach przedziału
czasu daje przechodniość; ciągłość wynika z lematu o sklejaniu.

\begin{figure}[!htbp]
\centering
\begin{tikzpicture}[>=Stealth,line cap=round,line join=round,scale=.9]
 \fill[manifoldblue!5,rounded corners=18pt]
  (-.4,-1.15) rectangle (10.2,2.45);
 \draw[manifoldblue!35,thick,rounded corners=18pt]
  (-.4,-1.15) rectangle (10.2,2.45);
 \coordinate (x) at (.5,.25);
 \coordinate (y) at (9.25,.25);
 \draw[manifoldred,very thick,-{Stealth[length=2.8mm]}]
  (x) .. controls (2.6,2.4) and (6.7,2.0) .. (y)
  node[pos=.48,above=4pt] {$\alpha$};
 \draw[manifoldgreen!80!black,very thick,-{Stealth[length=2.8mm]}]
  (x) .. controls (2.7,-1.0) and (6.5,-.7) .. (y)
  node[pos=.51,below=4pt] {$\beta$};
 \draw[manifoldgold!90!black,thick,-{Stealth[length=2.5mm]}]
  (4.9,1.5) -- (4.9,.35)
  node[pos=.52,right=4pt,font=\small] {deformacja};
 \fill[manifoldblue] (x) circle(2.4pt) (y) circle(2.4pt);
 \node[below left=3pt] at (x) {$x_0$};
 \node[below right=3pt] at (y) {$x_1$};
\end{tikzpicture}
\caption{Homotopia dróg przesuwa ich wnętrza,
ale zachowuje oba punkty końcowe.}
\label{fig:homotopia-drog}
\end{figure}

\begin{lemat}[Zmiana prędkości nie zmienia klasy]
\label{lem:reparametr-drogi}
Jeśli $r:I\to I$ jest ciągła, niemalejąca
i spełnia $r(0)=0$, $r(1)=1$, to
$\alpha\circ r$ jest homotopijna z $\alpha$
względem końców.
\end{lemat}
\begin{proof}
Połóżmy $r_t(s)=(1-t)r(s)+ts$.
Każde $r_t$ jest ciągłe, ma te same końce
i wartości w $I$. Funkcja
$H(s,t)=\alpha(r_t(s))$ jest ciągła.
Dla $t=0$ daje $\alpha\circ r$, dla $t=1$
daje $\alpha$, a końce pozostają stałe.
\end{proof}

\begin{twierdzenie}[Grupa podstawowa]
\label{thm:grupa-podstawowa}
Klasy homotopii pętli opartych w $x_0$,
z działaniem $[\alpha][\beta]=[\alpha*\beta]$,
tworzą grupę $\pi_1(X,x_0)$.
Element neutralny to stała pętla $c_{x_0}$,
a odwrotnością $[\alpha]$ jest $[\bar\alpha]$.
\end{twierdzenie}
\begin{proof}
Jeśli $H$ deformuje $\alpha$ w $\alpha'$,
a $K$ deformuje $\beta$ w $\beta'$
względem końców, sklejenie
\[
 L(s,t)=
 \begin{cases}
   H(2s,t),&s\leq\tfrac12,\\
   K(2s-1,t),&s\geq\tfrac12
 \end{cases}
\]
jest ciągłe, bo na wspólnej krawędzi
obie części mają wartość $x_0$.
Działanie nie zależy więc od reprezentantów.

Przy trzech pętlach $(\alpha*\beta)*\gamma$
przeznacza na $\alpha,\beta,\gamma$
odcinki czasu długości $1/4,1/4,1/2$,
zaś $\alpha*(\beta*\gamma)$ długości
$1/2,1/4,1/4$. Obie drogi przebiegają
te same kawałki w tej samej kolejności
i różnią się zmianą parametru.
Lemat~\ref{lem:reparametr-drogi} daje łączność
na klasach. Pętle $c_{x_0}*\alpha$
i $\alpha*c_{x_0}$ różnią się od $\alpha$
tylko chwilą postoju, więc
$[c_{x_0}]$ jest neutralny.

Dla odwrotności określmy
\[
 H(s,t)=
 \begin{cases}
  \alpha(2(1-t)s),&0\leq s\leq\tfrac12,\\
  \alpha(2(1-t)(1-s)),&\tfrac12\leq s\leq1.
 \end{cases}
\]
Przy $s=1/2$ obie gałęzie są równe
$\alpha(1-t)$, a przy $s=0,1$ mają
wartość $x_0$. Jest to homotopia
$\alpha*\bar\alpha$ do pętli stałej
względem punktu bazowego. Zamiana
$\alpha$ z $\bar\alpha$ daje też
$[\bar\alpha][\alpha]=[c_{x_0}]$.
\end{proof}

\begin{stwierdzenie}[Odwzorowania indukowane i punkt bazowy]
\label{prop:pi1-funktor}
Odwzorowanie ciągłe $f:(X,x_0)\to(Y,y_0)$
z $f(x_0)=y_0$ wyznacza homomorfizm
\[
 f_*:\pi_1(X,x_0)\to\pi_1(Y,y_0),
 \qquad [\alpha]\mapsto[f\circ\alpha].
\]
Zachodzi $(gf)_*=g_*f_*$ i
$(\id_X)_*=\id$. Homotopia odwzorowań
utrzymująca $x_0$ w $y_0$ daje
równe homomorfizmy. Jeśli
$u$ jest drogą od $x_0$ do $x_1$,
to $b_u([\alpha])=[\bar u*\alpha*u]$
jest izomorfizmem
$\pi_1(X,x_0)\to\pi_1(X,x_1)$.
\end{stwierdzenie}
\begin{proof}
Złożenie $f\circ H$ przenosi homotopię
pętli w homotopię ich obrazów, a
$f\circ(\alpha*\beta)
=(f\circ\alpha)*(f\circ\beta)$.
Równości dla złożeń wynikają stąd wprost.
Jeśli $F$ jest homotopią odwzorowań utrzymującą
punkt bazowy, $F(\alpha(s),t)$ jest
homotopią pętli opartych w $y_0$.
Wreszcie $b_u$ zachowuje iloczyn:
w środku złożenia pojawia się
$u*\bar u$, homotopijna z pętlą stałą.
Odwrotność stanowi $b_{\bar u}$.
\end{proof}

\begin{wniosek}[Niezmienniczość grupy podstawowej]
\label{cor:pi1-homotopia}
Jeśli $f\simeq g:X\to Y$ przez homotopię
$H$, a $u(t)=H(x_0,t)$ prowadzi od
$f(x_0)$ do $g(x_0)$, to
$g_*=b_u\circ f_*$. Każda równoważność
homotopijna $f:X\to Y$ daje więc
izomorfizm grup podstawowych,
po uwzględnieniu zmiany punktu bazowego.
\end{wniosek}
\begin{proof}
Dla pętli $\alpha$ rozpatrzmy kwadrat
$(s,t)\mapsto H(\alpha(s),t)$.
Jego dolny i górny bok to $f\alpha$
i $g\alpha$, a oba pionowe boki to
$u$. Brzeg kwadratu jest homotopijny do pętli stałej względem
wybranego narożnika: w wypukłym kwadracie ściągamy jego parametryzację
liniowo do tego narożnika, a następnie składamy z $H(\alpha(s),t)$.
Zatem w grupie z punktem
bazowym $g(x_0)$ mamy
$[g\alpha]=[\bar u*(f\alpha)*u]$.
To daje pierwszą równość.

Niech teraz $g$ będzie homotopijną
odwrotnością $f$ i oznaczmy
$y_0=f(x_0)$. Homotopia
$gf\simeq\id_X$ daje drogę $u$
od $gf(x_0)$ do $x_0$ i
izomorfizm $b_u\circ(gf)_*$.
Homotopia $fg\simeq\id_Y$ daje pewną drogę $v$
od $fg(y_0)$ do $y_0$ z $b_v\circ(fg)_*=\id$.
Droga $fu$ ma te same końce, choć może nie być równa $v$.
Ponieważ $b_{fu}\circ b_v^{-1}$ jest izomorfizmem,
również $b_{fu}\circ(fg)_*$ jest izomorfizmem grupy
$\pi_1(Y,y_0)$ na siebie
(każda zmiana punktu bazowego
jest izomorfizmem).
Z naturalności konkatenacji
$f_*\circ b_u=b_{f u}\circ f_*$
na odpowiednich grupach.
Zatem dla $g'_*:=b_u\circ g_*$
oba złożenia $g'_*f_*$
oraz $f_*g'_*$ są izomorfizmami.
Pierwsze zapewnia iniektywność
$f_*$, drugie jego surjektywność.
\end{proof}

Silna retrakcja deformacyjna
z~\ref{def:retrakt-deformacyjny} daje
izomorfizm grup podstawowych w punktach
podprzestrzeni. Dla przestrzeni wypukłej
każdą pętlę $\alpha$ można deformować
wzorem $(1-t)\alpha(s)+tx_0$ do pętli
stałej, nie ruszając punktu bazowego.

\section{Nakrycia: lokalnie wiele kopii jednej przestrzeni}
\label{sec:nakrycia}

Nakrycie przypomina nałożenie kilku arkuszy
na tę samą małą okolicę. Gdy idziemy po bazie,
wybranie początkowego arkusza jednoznacznie
ustala przebieg drogi na nakryciu.
Po zamkniętej drodze możemy wrócić na inny arkusz.

\begin{definicja}[Nakrycie i arkusze]
\index{nakrycie i arkusze@Nakrycie i arkusze}
\label{def:nakrycie}
Surjekcja ciągła $p:E\to X$ jest \emph{nakryciem},
jeśli każdy $x\in X$ ma otwarte otoczenie $U$
takie, że
\[
 p^{-1}(U)=\coprod_{\lambda\in\Lambda}V_\lambda,
 \qquad
 p|_{V_\lambda}:V_\lambda\xrightarrow{\cong}U,
\]
gdzie $V_\lambda$ są parami rozłącznymi
otwartymi podzbiorami $E$, a ograniczenia $p|_{V_\lambda}$
są homeomorfizmami.
Takie $U$ nazywamy \emph{równomiernie nakrytym},
a $V_\lambda$ \emph{arkuszami}.
Włókno $p^{-1}(x)$ jest dyskretne.
\end{definicja}

\begin{przyklad}[Prosta nad okręgiem]
\label{prz:exp-nakrycie}
Odwzorowanie
\[
 p:\R\to S^1,\qquad p(t)=e^{2\pi i t},
\]
jest nakryciem. Jeśli $U$ jest krótkim,
otwartym łukiem okręgu, wybierzmy przedział
$J$ długości mniejszej niż $1$, na którym
$p|_J$ jest homeomorfizmem na $U$.
Wtedy $p^{-1}(U)=\coprod_{k\in\Z}(J+k)$.
Włókno nad $1$ to $\Z$.
Podobnie $q_m:S^1\to S^1$, $q_m(z)=z^m$,
jest $m$-krotnym nakryciem dla $m\geq1$.
Dwukrotne nakrycie $\operatorname{Spin}(n)
\to\mathrm{SO}(n)$ pojawiło się już w
Twierdzeniu~\ref{thm:spin-nakrycie}.
\end{przyklad}

\begin{figure}[!htbp]
\centering
\begin{tikzpicture}[>=Stealth,line cap=round,line join=round,
 x={(1.55cm,0cm)},y={(.40cm,.37cm)},z={(0cm,1.25cm)}]
 % Okrąg u dołu, widoczna i zasłonięta połowa.
 \draw[manifoldblue!40,densely dashed,line width=1pt]
  plot[domain=0:180,samples=60,variable=\u]
   ({cos(\u)},{sin(\u)},0);
 \draw[manifoldblue,line width=1.8pt]
  plot[domain=180:360,samples=60,variable=\u]
   ({cos(\u)},{sin(\u)},0);
 % Włókno nad 1 oraz promień projekcji dla punktu t=1.25.
 \draw[manifoldblue!30,densely dotted]
  (1,0,0)--(1,0,2.9);
 \draw[manifoldgold!90!black,thick,dashed,
       -{Stealth[length=2.2mm]}]
  (0,1,1.65)--(0,1,0);
 \node[manifoldgold!70!black,right=4pt]
  at (0,1,.75) {$p$};
 % Fragment helisy: jeden obrót podnosi wysokość o jedną jednostkę.
 \draw[manifoldred,line width=2.1pt,
       -{Stealth[length=2.8mm]}]
  plot[domain=0:2.38,samples=180,variable=\t]
   ({cos(360*\t)},{sin(360*\t)},{.4+\t});
 \foreach \k in {0,1,2} {
  \fill[manifoldred] (1,0,{.4+\k}) circle(2.4pt);
  \node[manifoldred,right=5pt,font=\small]
   at (1,0,{.4+\k}) {$t=\k$};
 }
 \fill[manifoldblue] (1,0,0) circle(2.8pt);
 \node[manifoldblue,below right=3pt] at (1,0,0) {$1\in S^1$};
 \node[manifoldred!85!black,left=7pt]
  at (-1,0,2.7) {$\widetilde X\cong\R$};
 \node[manifoldblue!85!black,left=11pt]
  at (-1,0,-.22) {$X=S^1$};
\end{tikzpicture}
\caption{Nakrycie $p:\R\to S^1$, $p(t)=e^{2\pi i t}$,
przedstawione jako linia śrubowa nad okręgiem. Rzut pionowy
odpowiada $p$. Punkty linii śrubowej oznaczone $t=0,1,2$
leżą nad tym samym punktem $1\in S^1$.
Pokazano fragment linii śrubowej, która biegnie dalej
w obie strony.}
\label{fig:nakrycie-okragu}
\end{figure}

W poniższych dowodach użyjemy \emph{liczby Lebesgue'a}:
dla otwartego pokrycia zwartej przestrzeni metrycznej istnieje
$\delta>0$ takie, że każdy niepusty zbiór o średnicy mniejszej
niż $\delta$ leży w jednym elemencie pokrycia.
Oto uzasadnienie. Dla każdego $x$ wybierzmy $\varepsilon_x>0$
tak, aby $B(x,2\varepsilon_x)$ mieściła się w elemencie pokrycia.
Ze zwartości skończenie wiele kul $B(x_i,\varepsilon_{x_i})$
pokrywa przestrzeń. Wystarczy $\delta=\min_i\varepsilon_{x_i}$:
jeśli $y$ należy do jednej z tych mniejszych kul, wszystkie
punkty zbioru zawierającego $y$ i o średnicy $<\delta$
leżą w odpowiadającej jej większej kuli. Stosujemy ten fakt
do odcinka i kwadratu z ich zwykłą metryką.

\begin{twierdzenie}[Podnoszenie drogi i homotopii]
\label{thm:podnoszenie-nakrycie}
Niech $p:E\to X$ będzie nakryciem.
\begin{enumerate}[label=(\alph*),leftmargin=*]
\item Każda droga $\gamma:I\to X$ i wybór
      $e_0\in p^{-1}(\gamma(0))$ wyznaczają
      dokładnie jedną drogę
      $\widetilde\gamma:I\to E$ z
      $\widetilde\gamma(0)=e_0$ i
      $p\widetilde\gamma=\gamma$.
\item Jeśli $H:I\times I\to X$ jest
      homotopią dróg, a $\widetilde H(s,0)$
      jest zadanym podniesieniem $H(s,0)$,
      to istnieje dokładnie jedno ciągłe
      podniesienie $\widetilde H$ całego $H$.
      Gdy $H$ zachowuje końce dróg,
      podniesione końce również są stałe.
\end{enumerate}
\end{twierdzenie}
\begin{proof}
\emph{Droga.} Przeciwobrazy równomiernie
nakrytych otwartych $U$ tworzą otwarte
pokrycie odcinka $I$. Z liczby Lebesgue'a
tego pokrycia wybieramy podział
$0=t_0<\cdots<t_N=1$ tak drobny, aby
każdy obraz $\gamma([t_{j-1},t_j])$
mieścił się w jednym $U_j$.
Arkusz nad $U_1$ zawierający $e_0$
określa na pierwszym przedziale wzór
\[
 \widetilde\gamma(t)
 =(p|_{V_1})^{-1}(\gamma(t)).
\]
W punkcie $t_1$ wybieramy arkusz nad $U_2$
zawierający uzyskany punkt włókna i powtarzamy.
Wzory zgadzają się na końcach przedziałów,
więc sklejają się w drogę ciągłą.
Jeśli dwa podniesienia zgadzają się w pewnej
chwili, to lokalnie leżą w tym samym arkuszu
i są identyczne. Gdy w pewnej chwili
się różnią, leżą nad tym samym punktem
w różnych arkuszach, więc również w małej
okolicy chwili pozostają różne. Zbiór
chwil zgodności jest zatem otwarty
i domknięty w spójnym $I$ oraz
zawiera $0$, co daje jednoznaczność.

\emph{Homotopia.} Pokrycie $I^2$ zbiorami
$H^{-1}(U)$ ma liczbę Lebesgue'a.
Dzielimy kwadrat na skończoną siatkę
małych prostokątów, z których każdy ma
obraz w jednym równomiernie nakrytym $U$.
Na dolnej krawędzi mamy zadaną drogę.
Podnosimy prostokąty kolejno od dołu i
od lewej: na każdym wybieramy arkusz
zawierający już ustalony obraz jego narożnika
i stosujemy $(p|_V)^{-1}\circ H$.
Na wspólnej krawędzi oba wzory są
podniesieniami tej samej drogi i zgadzają
się w narożniku, więc na mocy (a) są równe.
Skończone sklejanie daje ciągłą
$\widetilde H$; lokalna jednoznaczność
na spójnym $I^2$ daje jej jedyność.
Jeśli $H(0,t)$ albo $H(1,t)$ jest stała,
jej podniesienie leży w dyskretnym włóknie,
więc też jest stałe.
\end{proof}

Stąd $p_*:\pi_1(E,e_0)\to\pi_1(X,p(e_0))$
jest iniektywny. Jeśli obraz pętli jest
homotopijny do stałej pętli względem bazy,
podnosimy homotopię od danej pętli w $E$.
Końcowa podniesiona pętla leży w dyskretnym
włóknie i jest stała.

\begin{twierdzenie}[Kryterium podnoszenia odwzorowań]
\label{thm:kryterium-podnoszenia}
Niech $Y$ będzie łukowo spójna i lokalnie
łukowo spójna, $f:Y\to X$ ciągła,
$p:E\to X$ nakryciem, a $p(e_0)=f(y_0)$.
Istnieje ciągła $\widetilde f:Y\to E$
z $p\widetilde f=f$ i
$\widetilde f(y_0)=e_0$ wtedy i tylko wtedy, gdy
\[
 f_*\pi_1(Y,y_0)\subseteq
 p_*\pi_1(E,e_0).
\]
Takie podniesienie jest jedyne.
\end{twierdzenie}
\begin{proof}
Gdy $\widetilde f$ istnieje,
$f=p\widetilde f$, zatem
$f_*=p_*\widetilde f_*$, co daje konieczność.
Załóżmy inkluzję. Dla $y\in Y$ wybierzmy
drogę $\gamma$ od $y_0$ do $y$ i
zdefiniujmy $\widetilde f(y)$ jako koniec
podniesienia $f\circ\gamma$ z początkiem $e_0$.
Gdy $\gamma'$ jest innym wyborem,
pętla $f\circ(\gamma'*\bar\gamma)$
ma klasę w obrazie $p_*$. Jest więc
homotopijna względem punktu bazowego
do obrazu pewnej pętli w $E$ o początku
$e_0$. Podnosząc tę homotopię za pomocą
Twierdzenia~\ref{thm:podnoszenie-nakrycie},
otrzymujemy, że podniesienie
$f\circ(\gamma'*\bar\gamma)$ wraca do $e_0$.
Druga część tego zamkniętego podniesienia jest podniesieniem
$f\circ\bar\gamma$ kończącym się w $e_0$.
Po odwróceniu jest, z jednoznaczności, podniesieniem
$f\circ\gamma$ zaczynającym się w $e_0$.
Podniesienia $f\circ\gamma$ i $f\circ\gamma'$ mają więc ten sam koniec.

Sprawdźmy ciągłość, której nie daje sama
definicja punkt po punkcie. Wybierzmy
równomiernie nakryte $U$ zawierające $f(y)$
oraz łukowo spójne otwarte $W\ni y$ z
$f(W)\subseteq U$. Niech $V$ będzie
arkuszem nad $U$ zawierającym $\widetilde f(y)$.
Dla $z\in W$ połączmy $y$ z $z$ drogą
w $W$ i dopiszmy ją do drogi od $y_0$
do $y$. Definicja i podnoszenie drogi
pokazują wtedy, że
$\widetilde f(z)=(p|_V)^{-1}(f(z))$.
Na $W$ odwzorowanie jest więc ciągłe.
Jedyność wynika z jedyności podniesienia
każdej drogi od $y_0$.
\end{proof}

\begin{twierdzenie}[Grupa podstawowa okręgu]
\label{thm:pi1-okragu}
Podnoszenie względem $p:\R\to S^1$
z Przykładu~\ref{prz:exp-nakrycie} daje
izomorfizm
\[
 \operatorname{wind}:\pi_1(S^1,1)\xrightarrow{\ \cong\ }\Z,
 \qquad [\gamma]\longmapsto
 \widetilde\gamma(1),
 \quad \widetilde\gamma(0)=0.
\]
Liczba po prawej to \emph{liczba owinięć}.
\end{twierdzenie}
\begin{proof}
Ponieważ $\gamma(1)=1$, koniec
podniesienia leży w $p^{-1}(1)=\Z$.
Homotopijne pętle dają homotopię
podniesień z nieruchomymi końcami,
więc liczba zależy tylko od klasy.
Jeśli podniesienie $\alpha$ kończy się w $k$,
to podniesienie następnej pętli $\beta$
zaczynające się w $k$ jest przesunięciem
o $k$ podniesienia zaczynającego się w $0$.
Stąd $\operatorname{wind}([\alpha*\beta])
=\operatorname{wind}([\alpha])
+\operatorname{wind}([\beta])$.
Pętla $\gamma_k(s)=e^{2\pi i ks}$
ma podniesienie $ks$, więc homomorfizm
jest surjektywny. Gdy podniesienie
$\widetilde\gamma$ wraca do $0$,
wzór $(s,t)\mapsto(1-t)\widetilde\gamma(s)$
deformuje ją w $\R$ do stałej drogi
z zachowaniem końców. Projekcja tej
homotopii do $S^1$ deformuje $\gamma$
do pętli stałej. Jądro jest zerowe.
\end{proof}

Z Przykładu~\ref{prz:retrakcja-pierscien}
otrzymujemy $\pi_1(\R^2\setminus\{0\})\cong\Z$.
W szczególności okrąg i punkt nie są
homotopijnie równoważne. Obliczenie
$\pi_1(S^{n})=0$ dla $n\geq2$
wyniknie niżej z twierdzenia van Kampena.

\begin{definicja}[Lokalna spójność i warunek półlokalny]
\index{lokalna spojnosc i warunek pollokalny@Lokalna spójność i warunek półlokalny}
\label{def:semilokalnie-prosto-spojna}
Przestrzeń $X$ jest \emph{lokalnie łukowo spójna},
gdy każdy punkt ma bazę otoczeń otwartych,
które są łukowo spójne. Jest
\emph{półlokalnie jednospójna}, gdy każdy
$x\in X$ ma otoczenie $U$, w którym
każda pętla oparta w $x$ jest homotopijna do stałej
względem punktu bazowego \emph{w całym $X$}. Nie wymagamy, aby
ta homotopia mieściła się w $U$.
Przestrzeń łukowo spójna z
$\pi_1(X,x_0)=0$ jest \emph{jednospójna}.
\end{definicja}

Każda rozmaitość topologiczna spełnia
lokalne warunki z tej definicji, gdyż
w otwartej mapie współrzędnych można
wybrać małą kulę. Warunek półlokalny
jest ważny: bez niego nakrycie uniwersalne
może w ogóle nie istnieć.

\begin{twierdzenie}[Konstrukcja nakrycia uniwersalnego]
\label{thm:uniwersalne-nakrycie}
Jeśli $X$ jest łukowo spójna, lokalnie
łukowo spójna i półlokalnie jednospójna,
to ma nakrycie $q:\widetilde X\to X$
z jednospójną przestrzenią $\widetilde X$.
\end{twierdzenie}
\begin{proof}
Ustalmy $x_0\in X$. Punktem $\widetilde X$
niech będzie klasa $[\gamma]$ drogi w $X$
zaczynającej się w $x_0$, przy czym
homotopie zachowują oba końce.
Definiujemy $q([\gamma])=\gamma(1)$.
Istnieje baza otwartych łukowo spójnych
zbiorów $U\subseteq X$, w których
każda pętla jest zerowa po włączeniu do $X$:
najpierw wybieramy otoczenie z warunku
półlokalnego, następnie w nim małe
łukowo spójne otoczenie z lokalnej
łukowej spójności. Każdą pętlę opartą w innym punkcie
takiego $U$ przenosimy do jego wyróżnionego punktu drogą w $U$;
izomorfizm zmiany punktu bazowego pokazuje, że ona też jest zerowa w $X$.

Dla drogi $\gamma$ kończącej się w $x\in U$
niech $W([\gamma],U)$ składa się z klas
$[\gamma*\eta]$, gdzie $\eta$ jest drogą
w $U$ od $x$ do dowolnego $y\in U$.
Jeśli dwie takie drogi $\eta,\eta'$
mają ten sam koniec, to
$\eta*\bar\eta'$ jest pętlą w $U$,
zerową w $X$. Zatem $[\gamma*\eta]
=[\gamma*\eta']$. Otrzymujemy bijekcję
$q:W([\gamma],U)\to U$. Bierzemy wszystkie
te zbiory jako bazę topologii. Sprawdźmy warunek przecięć:
jeśli $[\beta]\in W([\gamma],U)\cap W([\gamma'],U')$,
wybieramy otwarte łukowo spójne $O$ z naszej bazy,
zawierające $\beta(1)$ i zawarte w $U\cap U'$.
Wtedy $W([\beta],O)$ zawiera $[\beta]$ i mieści się w obu zbiorach:
do drogi kończącej się w $\beta(1)$ dopisujemy tylko drogę w $O$.
To dowodzi, że rzeczywiście otrzymujemy bazę.
Ograniczenie $q$ do każdego $W$ jest ciągłe, ponieważ przeciwobrazy
otwartych podzbiorów $U$ są sumami takich mniejszych zbiorów bazowych;
jest też otwarte, bo mniejsze $W([\beta],O)$ przechodzą na otwarte $O$.
Wcześniej wykazana bijekcja jest więc homeomorfizmem.
Jeśli dwa $W$ nad tym samym
$U$ mają wspólny punkt, ich drogi początkowe
różnią się drogą w $U$ z dokładnością
do homotopii w $X$, więc zbiory $W$
są równe. Dla mniejszych otoczeń
używamy ograniczenia tych samych odwzorowań.
Stąd $q^{-1}(U)$ jest rozłączną sumą
arkuszy $W([\gamma],U)$: $q$ jest nakryciem.

Jest ono łukowo spójne. Drogę od klasy
stałej $[c_{x_0}]$ do $[\gamma]$
otrzymujemy, biorąc po kolei klasy
odcinków $\gamma|_{[0,t]}$ po
przeparametryzowaniu na $I$; lokalnie
leżą one w jednym $W$, więc tworzą
ciągłą drogę. Jeśli pętla
$\lambda$ w $\widetilde X$ jest oparta
w $[c_{x_0}]$, jej rzut $\gamma=q\lambda$
ma podniesienie $\lambda$ i kończy
się w tej samej klasie drogi.
Konstrukcja podniesienia pokazuje,
że ta klasa jest $[\gamma]$, więc
$\gamma$ jest homotopijna do stałej
pętli względem bazowego końca.
Podnosimy tę homotopię, zaczynając
od $\lambda$, za pomocą
Twierdzenia~\ref{thm:podnoszenie-nakrycie}.
Na końcu dostajemy stałą pętlę
w dyskretnym włóknie. Zatem
$\pi_1(\widetilde X)=0$.
\end{proof}

\begin{uwaga}[Skąd podgrupy grupy podstawowej?]
\label{uwaga:nakrycia-podgrupy}
Dla spójnego nakrycia $p:(E,e_0)\to(X,x_0)$
iniekcja $p_*$ wyróżnia podgrupę
$H=p_*\pi_1(E,e_0)\leq\pi_1(X,x_0)$.
W warunkach Twierdzenia~\ref{thm:uniwersalne-nakrycie}
każda taka podgrupa daje nakrycie:
na klasach dróg w $\widetilde X$ grupa
$\pi_1(X,x_0)$ działa przez
$[\omega]\cdot[\gamma]=[\omega*\gamma]$,
a iloraz $H\backslash\widetilde X$
jest nakryciem $X$. Istotnie, działanie przenosi
$W([\gamma],U)$ na $W([\omega*\gamma],U)$,
więc jest działaniem przez homeomorfizmy. Jest wolne:
z $[\omega*\gamma]=[\gamma]$ po dopisaniu $\bar\gamma$
wynika $[\omega]=1$.
Co więcej, jeśli $hW\cap W\ne\varnothing$, to $hW=W$,
bo są to arkusze nad tym samym $U$. Ponieważ $q$ jest
iniektywne na $W$ i $q(hz)=q(z)$, element $h$ ustala punkt $z\in W$;
stąd $h=1$. Różne przesunięcia jednego arkusza są więc rozłączne.
Projekcja na orbity jest otwarta: przeciwobraz obrazu zbioru
otwartego $O$ jest sumą otwartych zbiorów $hO$.
Obrazy arkuszy w ilorazie są zatem otwarte i każdy odwzorowuje się
homeomorficznie na $U$. To uzasadnia własność nakrycia.
Przestrzeń ilorazowa jest łukowo spójna jako obraz $\widetilde X$.
Podniesienie pętli $\gamma$ z $x_0$ kończy się w orbicie $[\gamma]$;
wraca do orbity drogi stałej dokładnie wtedy, gdy $[\gamma]\in H$.
Ponieważ obraz homomorfizmu indukowanego przez nakrycie składa się
właśnie z klas mających zamknięte podniesienie, skonstruowane nakrycie
realizuje dokładnie podgrupę $H$.
Z drugiej strony kryterium
\ref{thm:kryterium-podnoszenia} podnosi
$q:\widetilde X\to X$ do $E$, z ustaloną wartością $e_0$.
Przy obecnych założeniach $E$ jest lokalnie łukowo spójna,
więc jej spójność daje łukową spójność: składowe łukowe są otwarte.
Każdy punkt $E$ można połączyć z $e_0$ drogą; rzut tej drogi
pokazuje, że podniesione odwzorowanie $\widetilde X\to E$ jest surjektywne.
Dwa punkty $[\gamma],[\gamma']$ nad tym samym punktem $X$
trafiają w ten sam punkt $E$
dokładnie wtedy, gdy $[\gamma'*\bar\gamma]\in H$.
W jedną stronę składamy podniesione drogi o wspólnym końcu,
otrzymując pętlę w $E$; w drugą używamy podnoszenia homotopii
jak w dowodzie kryterium podnoszenia. Ten warunek jest też
równoważny przynależności do jednej orbity, gdyż
$[\gamma'*\bar\gamma]\cdot[\gamma]=[\gamma']$.
Lokalnie obie projekcje są
homeomorfizmami na $U$, co daje
$E\cong H\backslash\widetilde X$
nad $X$. Zmiana wybranego punktu
$e_0$ sprzęga podgrupę $H$: droga w $E$ między dwoma punktami
nad $x_0$ rzutuje się na pętlę $u$, a zmiana bazy zastępuje
$H$ podgrupą $[\bar u]H[u]$.
Odwrotnie, dowolną pętlę $u$ w $X$ można podnieść z $e_0$;
jej koniec wybiera punkt włókna realizujący właśnie to sprzężenie.
Tak powstaje klasyfikacja spójnych
nakryć przez klasy sprzężoności
podgrup $\pi_1(X,x_0)$.
\end{uwaga}

\section{Twierdzenie van Kampena i pierwsze obliczenia}
\label{sec:van-kampen}

Nakrycie okręgu mierzy jedną liczbę owinięć.
Gdy przestrzeń składa się z kilku części,
potrzebujemy sposobu łączenia informacji
o pętlach w tych częściach.

\begin{twierdzenie}[van Kampen: dwa otwarte zbiory]
\label{thm:van-kampen}
Niech $X=U\cup V$, gdzie $U,V$ są otwarte,
łukowo spójne, a $U\cap V$ jest łukowo
spójne i zawiera $x_0$. Jeśli
$G_U=\pi_1(U,x_0)$, $G_V=\pi_1(V,x_0)$
oraz $G_0=\pi_1(U\cap V,x_0)$, to
\[
 \pi_1(X,x_0)\cong
 (G_U*G_V)/
 \left\langle\!\left\langle
 (i_U)_*(h)(i_V)_*(h)^{-1}:h\in G_0
 \right\rangle\!\right\rangle.
\]
Iloczyn wolny i redukcję słów wyjaśnia
Definicja~\ref{def:free-groups-5}, a zapis
$\langle\!\langle\cdot\rangle\!\rangle$
oznacza domknięcie normalne z
Definicji~\ref{def:normal-quotient-5}.
\end{twierdzenie}
\begin{proof}
Włączenia $U,V\hookrightarrow X$
określają homomorfizm
$\Phi:G_U*G_V\to\pi_1(X,x_0)$.
Pętla w $U\cap V$ daje ten sam obraz
niezależnie od wyboru strony.
Zatem wszystkie relacje z tezy
leżą w $\ker\Phi$.

\emph{Surjektywność.} Dla pętli
$\alpha:I\to X$ przeciwobrazy $U,V$
pokrywają zwarty odcinek. Dzielimy
go na tyle małych przedziałów,
aby obraz każdego leżał w $U$ albo $V$.
Sąsiednie kawałki przypisane temu
samemu zbiorowi łączymy. Każdy
punkt przejścia między kawałkiem
w $U$ i kawałkiem w $V$ leży w
$U\cap V$. Łączymy go z $x_0$
pomocniczą drogą w $U\cap V$.
Jeśli $\lambda_j$ prowadzi od $x_0$
do $j$-tego punktu przejścia,
to z kawałka $\alpha_j$ budujemy
pętlę $\lambda_{j-1}*\alpha_j*
\bar\lambda_j$ całkowicie w $U$
albo w $V$. Iloczyn tych pętli
odtwarza $\alpha$: sąsiednie
$\bar\lambda_j*\lambda_j$
kasują się na mocy
Twierdzenia~\ref{thm:grupa-podstawowa}.

Zmiana drogi $\lambda_j$ na $\lambda'_j$
wyznacza pętlę $h=\lambda'_j*\bar\lambda_j$ w $U\cap V$
\emph{opartą w $x_0$}. Poprzedni czynnik słowa mnożymy z prawej
przez $h^{-1}$, a następny z lewej przez $h$.
Oba zapisy $h$ są równe w ilorazie z tezy, więc się skracają.
Również dopisanie podziału wewnątrz
jednego otwartego zbioru nie zmienia
otrzymanej klasy słowa.

\emph{Jądro.} Przypuśćmy, że słowo
z $G_U*G_V$ daje pętlę homotopijną
do stałej w $X$. Wypełnijmy ją
homotopią $F:I^2\to X$, stałą na
pozostałych trzech bokach.
Zwartość kwadratu pozwala wybrać
siatkę tak drobną, by obraz
każdego małego prostokąta leżał
w całości w $U$ lub w $V$.
Podział dolnego boku dobieramy także tak, aby zawierał końce
pętli reprezentujących litery wyjściowego słowa.
Dla każdego wierzchołka $z$ siatki wybierzmy drogę $\mu_z$
od $x_0$ do $F(z)$. Jeśli $F(z)\in U\cap V$, wybieramy ją
w przecięciu; jeśli $F(z)$ należy tylko do jednego zbioru,
wybieramy drogę w tym zbiorze. Jest to możliwe dzięki założeniom
o łukowej spójności. Przy $F(z)=x_0$ wybieramy drogę stałą.

Zorientowanej krawędzi $e:z\to w$ przypisujemy klasę
$[\mu_z*(F|_e)*\bar\mu_w]$ w grupie zbioru zawierającego
obraz sąsiedniego prostokąta. Obie drogi pomocnicze leżą w tym
zbiorze. Jeżeli dwa sąsiednie prostokąty mają różne oznaczenia,
ich wspólna krawędź i obie drogi pomocnicze leżą w $U\cap V$,
więc przypisania dają ten sam element ilorazu z tezy.
Odwrócenie krawędzi daje element odwrotny.

Iloczyn czterech elementów wokół jednego prostokąta jest jedynką:
drogi pomocnicze skracają się parami, a jego brzeg jest homotopijny
do stałej we właściwym $U$ lub $V$, przez mapę $F$ na tym prostokącie.
Możemy zatem zastępować dwa kolejne boki prostokąta pozostałymi
dwoma, z tymi samymi końcami, bez zmiany elementu ilorazu.
Przesuwając tak łamaną po kolejnych prostokątach, wiersz po wierszu,
zamieniamy dolny bok kwadratu na drogę wzdłuż pozostałych trzech boków.
Te trzy boki są stałe. Na dolnym boku drogi pomocnicze wewnątrz
każdej pierwotnej litery skracają się w jej grupie, a przy końcach
liter są stałe. Odtwarzamy więc dokładnie wyjściowe słowo,
które okazało się jedynką w ilorazie. Każdy krok używał tylko
relacji wewnątrz obu grup i relacji pochodzących z przecięcia.
Stąd $\ker\Phi$ zawiera się w
wskazanym domknięciu normalnym relacji, co kończy dowód.
\end{proof}

\begin{przyklad}[Dwa okręgi złączone w punkcie]
\label{prz:pi1-dwa-okregi}
Niech $X=S^1_a\vee S^1_b$.
Wybierzmy na każdym okręgu po jednym
punkcie różnym od punktu sklejenia
$x_0$. Usuwając punkt z drugiego
okręgu, otrzymujemy otwarty $U$
retrakcyjny na $S^1_a$; analogicznie
$V$ retrakcyjny na $S^1_b$.
Przecięcie $U\cap V$ jest ściągalne
(to dwa otwarte łuki spotykające się
w $x_0$). Wobec tego
\[
 \pi_1(S^1\vee S^1,x_0)
 \cong\Z*\Z=:F(a,b).
\]
Słowa $ab$ i $ba$ są różne:
grupa podstawowa nie musi być przemienna.
\end{przyklad}

\begin{przyklad}[Sfera wymiaru co najmniej dwa]
\label{prz:pi1-sfera}
Niech $n\geq2$. Zbiory
$U=S^n\setminus\{N\}$ i
$V=S^n\setminus\{S\}$, gdzie $N,S$
są biegunami, są homeomorficzne
z $\R^n$, więc mają trywialne
grupy podstawowe. Ich przecięcie
jest łukowo spójne: deformuje się
na równik $S^{n-1}$, który jest
łukowo spójny dla $n-1\geq1$.
Twierdzenie~\ref{thm:van-kampen}
daje $\pi_1(S^n)=0$.
Stąd $\pi_1(\R^{n+1}\setminus\{0\})=0$
dla $n\geq2$, co odróżnia ten
przypadek od płaszczyzny.
\end{przyklad}

\section{Sympleksy i triangulacje}
\label{sec:sympleksy}

Sympleks jest najprostszym wielowymiarowym
klockiem: odcinek w wymiarze $1$, trójkąt
w wymiarze $2$, czworościan w wymiarze $3$.
Współrzędne barycentryczne opisują
zarazem punkt i to, na której ścianie leży.

\begin{definicja}[Sympleks standardowy i ściany]
\index{sympleks standardowy i sciany@Sympleks standardowy i ściany}
\label{def:sympleks-standardowy}
Dla $n\geq0$ \emph{standardowy $n$-sympleks}
to
\[
 \Delta^n=
 \left\{(t_0,\ldots,t_n)\in\R^{n+1}:
 t_i\geq0,\ \sum_{i=0}^{n}t_i=1\right\}.
\]
Jego wierzchołki to wektory $e_0,\ldots,e_n$.
\emph{Wnętrze} oznacza $t_i>0$ dla wszystkich $i$
(względem afinicznej otoczki sympleksu).
Dla $n\geq1$ \emph{ściana kowymiaru jeden} naprzeciw $e_i$
jest zadana równaniem $t_i=0$ i jest kopią $\Delta^{n-1}$.
Ogólniej ścianę wyznacza dowolny niepusty podzbiór wierzchołków;
jest właściwa, gdy pomija choć jeden z nich.
Suma wszystkich właściwych ścian to $\partial\Delta^n$.
Przy $n=0$ nie ma niepustych właściwych ścian i $\partial\Delta^0=\varnothing$.
Punkty $v_0,\ldots,v_n\in\R^N$
są \emph{afinicznie niezależne}, jeśli
$v_1-v_0,\ldots,v_n-v_0$ są liniowo
niezależne. Ich otoczka wypukła
$[v_0,\ldots,v_n]$ jest
\emph{geometrycznym $n$-sympleksem}.
\end{definicja}

Każdy punkt sympleksu ma postać
$x=\sum_i t_iv_i$. Współczynniki
$t_i$ są jednoznaczne: jeśli
$\sum_i(t_i-s_i)v_i=0$ i obie sumy
współczynników są równe $1$, to
$\sum_{i=1}^n(t_i-s_i)(v_i-v_0)=0$;
afiniczna niezależność wymusza
$t_i=s_i$ dla wszystkich $i$.
W szczególności usunięcie jednego
wierzchołka daje ścianę
o wymiarze $n-1$.

\begin{figure}[!htbp]
\centering
\begin{tikzpicture}[line cap=round,line join=round]
 \node[font=\bfseries] at (0,3.45) {$\Delta^0$};
 \node[font=\bfseries] at (4.05,3.45) {$\Delta^1$};
 \node[font=\bfseries] at (8.63,3.45) {$\Delta^2$};
 % Wymiar zero.
 \fill[manifoldblue!10] (0,1.15) circle(.26);
 \fill[manifoldred] (0,1.15) circle(3.2pt);
 \node[below=9pt] at (0,1.15) {$v_0$};
 % Wymiar jeden.
 \coordinate (b0) at (2.75,1.15);
 \coordinate (b1) at (5.35,1.15);
 \draw[manifoldblue,line width=2.6pt] (b0)--(b1);
 \fill[manifoldred] (b0) circle(3.2pt)
  (b1) circle(3.2pt);
 \node[below=8pt] at (b0) {$v_0$};
 \node[below=8pt] at (b1) {$v_1$};
 % Wymiar dwa.
 \coordinate (c0) at (7.00,.35);
 \coordinate (c1) at (10.28,.35);
 \coordinate (c2) at (8.21,2.45);
 \fill[manifoldblue!13] (c0)--(c1)--(c2)--cycle;
 \draw[manifoldblue,line width=2pt] (c0)--(c1)--(c2)--cycle;
 \fill[manifoldred] (c0) circle(3.2pt)
  (c1) circle(3.2pt) (c2) circle(3.2pt);
 \node[below left=3pt] at (c0) {$v_0$};
 \node[below right=3pt] at (c1) {$v_1$};
 \node[above=4pt] at (c2) {$v_2$};
\end{tikzpicture}
\caption{Sympleksy wymiarów $0,1,2$: punkt,
odcinek i wypełniony trójkąt. Wymiar
$n$ odpowiada $n+1$ afinicznie niezależnym
wierzchołkom; ściany powstają przez ich usunięcie.}
\label{fig:sympleks-trojkat}
\end{figure}

\begin{stwierdzenie}[Brzeg sympleksu jest sferą]
\label{prop:brzeg-sympleksu}
Dla $n\geq1$ przestrzeń $\Delta^n$
jest homeomorficzna z dyskiem $D^n$,
a $\partial\Delta^n$ z $S^{n-1}$.
\end{stwierdzenie}
\begin{proof}
Pracujemy w $n$-wymiarowej płaszczyźnie
afinicznej $P=\{\sum t_i=1\}$.
Środek $c=(1/(n+1),\ldots,1/(n+1))$
leży we wnętrzu. Każdy promień z $c$
w $P$ przecina brzeg dokładnie raz:
nierówności $t_i\geq0$ wyznaczają
na nim odcinek od $c$ do pierwszej
ściany, na której pewna współrzędna
staje się zerem.
Przypisanie punktowi brzegowemu
$x$ kierunku $(x-c)/\|x-c\|$
jest ciągłą bijekcją
$\partial\Delta^n\to S^{n-1}$
po utożsamieniu $P-c\cong\R^n$.
Dziedzina jest zwarta, a sfera
Hausdorffa, więc odwrotność jest
ciągła. Jeśli $b(u)$ jest tą odwrotnością dla kierunku
$u\in S^{n-1}$, wzór $ru\mapsto c+r(b(u)-c)$, $0<r\leq1$,
z wartością $c$ w zerze daje ciągłą bijekcję z $D^n$ na $\Delta^n$.
Ciągłość w zerze wynika z ograniczoności $b(u)-c$.
Ponownie zwartość i własność Hausdorffa dają homeomorfizm,
który na brzegu jest poprzednią identyfikacją. Dla $n=1$ brzeg
składa się z dwóch punktów,
czyli jest $S^0$.
\end{proof}

\begin{definicja}[Skończony kompleks symplicjalny]
\index{skonczony kompleks symplicjalny@Skończony kompleks symplicjalny}
\label{def:kompleks-symplicjalny}
Niech $V$ będzie skończonym zbiorem
wierzchołków. \emph{Kompleks symplicjalny}
$K$ to rodzina niepustych skończonych
podzbiorów $V$ taka, że jeśli
$\sigma\in K$ i $\varnothing\ne\tau\subseteq\sigma$,
to $\tau\in K$. Zbiór o $n+1$
wierzchołkach nazywamy
\emph{$n$-sympleksem} $K$.
Jego \emph{realizacja geometryczna} to
\[
 |K|=\left\{\sum_{v\in V}t_v e_v:
 t_v\geq0,\ \sum_v t_v=1,\
 \{v:t_v>0\}\in K\right\}\subseteq\R^{V}.
\]
Kompleks $L\subseteq K$ zamknięty
na przechodzenie do ścian jest
\emph{podkompleksem}.
\emph{Skończona triangulacja} przestrzeni $X$
to homeomorfizm $|K|\to X$. W tej sekcji ograniczamy się do
skończonych kompleksów; ich realizacje są zwarte. Triangulacje
przestrzeni niezwartych wymagają dopuszczenia nieskończenie wielu sympleksów.
\end{definicja}

W tej definicji dwie różne ściany
mają dokładnie to wspólne, co
wynika ze wspólnych wierzchołków.
Na przykład trzy krawędzie
$[v_0,v_1]$, $[v_1,v_2]$,
$[v_2,v_0]$, bez wnętrza trójkąta,
realizują okrąg. Po dodaniu
$[v_0,v_1,v_2]$ otrzymujemy dysk.
Słowo \emph{podkompleks} jest
ważniejsze niż dowolny podzbiór:
zawiera wraz z sympleksem
wszystkie jego ściany.

\begin{uwaga}[Nie każdy podział na trójkąty
jest kompleksem symplicjalnym]
\label{uwaga:delta-motywacja}
Kwadrat z utożsamionymi przeciwległymi
bokami przedstawia torus. Gdy podzielimy
go jedną przekątną, po sklejeniu wszystkie
cztery narożniki stają się jednym
wierzchołkiem, a każda z dwóch
trójkątnych części nadal ma sens.
Nie jest to \emph{kompleks symplicjalny}
z jednym wierzchołkiem: trójkąt
potrzebuje trzech różnych wierzchołków.
Wygodniejszym opisem są
$\Delta$-kompleksy.
\end{uwaga}

\begin{definicja}[$\Delta$-kompleks]
\label{def:delta-kompleks}
\emph{$\Delta$-kompleks} $X$ jest
przestrzenią Hausdorffa z rodziną ciągłych odwzorowań charakterystycznych
$\sigma_\alpha:\Delta^{n_\alpha}\to X$:
\begin{enumerate}[label=(\roman*),leftmargin=*]
\item ograniczenie $\sigma_\alpha$ do wnętrza
      jest iniektywne, a obrazy wnętrz
      różnych sympleksów są rozłączne
      i pokrywają $X$;
\item ograniczenie $\sigma_\alpha$ do każdej
      ściany jest jednym z odwzorowań
      charakterystycznych niższego wymiaru
      po standardowym afinicznym
      utożsamieniu tej ściany z sympleksem, zachowującym
      kolejność pozostałych wierzchołków;
\item $A\subseteq X$ jest domknięty dokładnie
      wtedy, gdy $\sigma_\alpha^{-1}(A)$
      jest domknięty w każdym
      $\Delta^{n_\alpha}$.
\end{enumerate}
Różne ściany tego samego sympleksu
mogą być tu sklejone, o ile zgodne
są odwzorowania na ich dalszych ścianach.
\end{definicja}

Ostatni warunek to \emph{topologia słaba}
względem odwzorowań charakterystycznych:
odwzorowanie z $X$ do innej przestrzeni
jest ciągłe dokładnie wtedy, gdy jego
złożenia z każdą $\sigma_\alpha$ są ciągłe.
Jest to topologia ilorazowa z
Definicji~\ref{def:topologia-ilorazowa}
na rozłącznej sumie wszystkich sympleksów.

\begin{przyklad}[Okrąg i torus jako $\Delta$-kompleksy]
\label{prz:delta-torus}
Okrąg ma jeden $0$-sympleks $v$
i jeden $1$-sympleks, którego oba końce
przechodzą w $v$. Dla torusa
podzielmy kwadrat $[0,1]^2$
przekątną i sklejmy
$(0,t)\sim(1,t)$ oraz
$(s,0)\sim(s,1)$. Otrzymujemy
jeden wierzchołek, trzy krawędzie
(poziomą $a$, pionową $b$ i przekątną
$c$) oraz dwa $2$-sympleksy. Porządkujemy ich wierzchołki jako
$((0,0),(0,1),(1,1))$ i $((0,0),(1,0),(1,1))$.
Oba porządki nadają krawędziom poziomym kierunek w prawo,
pionowym w górę, a przekątnej kierunek od $(0,0)$ do $(1,1)$.
Mapy ścian są zatem zgodne także z porządkiem wierzchołków.
Iloraz jest homeomorficzny z
$(\R/\Z)\times(\R/\Z)$:
odwzorowanie $(s,t)\mapsto([s],[t])$
identyfikuje dokładnie wskazane
punkty brzegowe, jest ciągłe
i przechodzi na bijekcję
między przestrzenią zwartą
a przestrzenią Hausdorffa.
\end{przyklad}

\begin{figure}[!htbp]
\centering
\begin{tikzpicture}[>=Stealth,line cap=round,line join=round,
 x=1cm,y=1cm]
 \coordinate (a) at (0,0);
 \coordinate (b) at (4.4,0);
 \coordinate (c) at (4.4,4.4);
 \coordinate (d) at (0,4.4);
 \fill[manifoldblue!8] (a)--(b)--(c)--cycle;
 \fill[manifoldgreen!10] (a)--(d)--(c)--cycle;
 \draw[manifoldblue,very thick] (a)--(b)--(c)--(d)--cycle;
 \draw[manifoldgold!90!black,thick] (a)--(c);
 \draw[manifoldgold!90!black,thick,-{Stealth[length=2.7mm]}]
  (1.8,1.8)--(2.7,2.7);
 \draw[manifoldred,very thick,-{Stealth[length=2.7mm]}]
  (1.3,0)--(2.6,0);
 \draw[manifoldred,very thick,-{Stealth[length=2.7mm]}]
  (1.3,4.4)--(2.6,4.4);
 \draw[manifoldgreen!75!black,very thick,-{Stealth[length=2.7mm]}]
  (0,1.3)--(0,2.6);
 \draw[manifoldgreen!75!black,very thick,-{Stealth[length=2.7mm]}]
  (4.4,1.3)--(4.4,2.6);
 \node[manifoldred,below=2pt] at (3.2,0) {$a$};
 \node[manifoldgreen!65!black,right=2pt] at (4.4,3.2) {$b$};
 \node[manifoldgold!65!black,above left=2pt]
  at (2.6,2.6) {$c$};
 \foreach \p in {a,b,c,d}
  \fill[manifoldblue] (\p) circle(2.3pt);
 \node[font=\small] at (1.42,2.48) {$\Delta^2_1$};
 \node[font=\small] at (3.17,1.74) {$\Delta^2_2$};
\end{tikzpicture}
\caption{Torus z dwóch sympleksów: boki o tych samych
kolorach i kierunkach sklejamy ze sobą. Przekątna
pozostaje trzecią krawędzią. Wszystkie cztery narożniki
reprezentują ten sam wierzchołek torusa.}
\label{fig:delta-torus}
\end{figure}

\section{CW-kompleksy i dołączanie komórek}
\label{sec:cw-kompleksy}

Trójkąty są wygodne do obliczeń, lecz
potrafią wymagać wielu wierzchołków.
Komórka pozwala dołączyć naraz cały
otwarty dysk, przy czym jego brzeg
może biec po starszej części przestrzeni
w sposób nietrywialny.

\begin{definicja}[Dołączenie komórki]
\index{dolaczenie komorki@Dołączenie komórki}
\label{def:dolaczenie-komorki}
Niech $Y$ będzie przestrzenią i
$\varphi:S^{n-1}=\partial D^n\to Y$
odwzorowaniem ciągłym, $n\geq1$.
\emph{Dołączeniem $n$-komórki} jest
przestrzeń ilorazowa
\[
 Y\cup_\varphi D^n
 :=(Y\sqcup D^n)/(z\sim\varphi(z)
 \text{ dla }z\in S^{n-1}).
\]
Obraz wnętrza dysku oznaczamy $e^n$.
Odwzorowanie z dysku do ilorazu jest
\emph{odwzorowaniem charakterystycznym}
tej komórki; $\varphi$ nazywamy
\emph{odwzorowaniem dołączającym}.
W wymiarze $0$ po prostu dołączamy
izolowany punkt, gdyż
$S^{-1}=\varnothing$.
\end{definicja}

Naturalne włączenie $Y$ w przestrzeń dołączoną jest osadzeniem
na zbiór domknięty. Dla każdego domkniętego $F\subseteq Y$
przeciwobraz jego obrazu w $Y\sqcup D^n$ jest równy
$F\sqcup\varphi^{-1}(F)$ i jest domknięty; dotyczy to również $F=Y$.
Ponieważ na punktach $Y$ nie wprowadzamy nowych utożsamień,
dostajemy zamknięte osadzenie. Wnętrze $D^n$ jest natomiast
otwartym zbiorem nasyconym, bez żadnych utożsamień, więc przechodzi
homeomorficznie na otwartą komórkę $e^n$.

Topologia ilorazowa z
Definicji~\ref{def:topologia-ilorazowa}
ma tu konkretne znaczenie. Jeśli
$F:Y\cup_\varphi D^n\to Z$
ma być ciągła, wystarczy podać
ciągłe $F_Y:Y\to Z$ i
$F_D:D^n\to Z$ zgodne na brzegu:
$F_D(z)=F_Y(\varphi(z))$.
Istotnie, takie odwzorowania dają ciągłe
odwzorowanie na rozłącznej sumie, stałe
na każdej klasie sklejenia;
uniwersalna własność ilorazu
przenosi je na iloraz.
To jest praktyczny sposób
konstruowania odwzorowań na przestrzeniach
sklejonych.

\begin{definicja}[CW-kompleks]
\index{cw-kompleks@CW-kompleks}
\label{def:cw-kompleks}
\emph{CW-kompleks} $X$ jest przestrzenią Hausdorffa budowaną
ze szkieletów $X^0\subseteq X^1
\subseteq X^2\subseteq\cdots$:
$X^0$ jest dyskretnym zbiorem punktów,
a $X^n$ powstaje z $X^{n-1}$
przez dołączenie rodziny $n$-dysków
odwzorowaniami $S^{n-1}\to X^{n-1}$.
Przestrzeń $X=\bigcup_n X^n$
ma \emph{topologię słabą}:
$A\subseteq X$ jest domknięty,
dokładnie wtedy, gdy przeciwobraz $A$ pod każdym
odwzorowaniem charakterystycznym $D^n\to X$
jest domknięty. Wymagamy ponadto
\emph{skończoności domknięć komórek}:
domknięcie każdej komórki
spotyka tylko skończenie wiele
komórek. Litery CW przypominają
te dwa warunki
(\emph{closure-finite}, \emph{weak topology}).
\emph{Podkompleks CW} to suma komórek
domknięta w $X$; jeśli zawiera komórkę,
zawiera także całe jej domknięcie.
\end{definicja}

\begin{figure}[!htbp]
\centering
\begin{tikzpicture}[>=Stealth,line cap=round,line join=round]
 \node[font=\bfseries] at (0,3.02) {$X^0$};
 \node[font=\bfseries] at (4.33,3.02) {$X^1$};
 \node[font=\bfseries] at (8.68,3.02) {$X^2$};
 % 0-komórka.
 \fill[manifoldred] (0,1.39) circle(3.7pt);
 \node[below=7pt] at (0,1.39) {$e^0=\{v\}$};
 % Dołączenie 1-komórki: oba końce trafiają do v.
 \draw[manifoldblue,line width=2.5pt]
  (4.33,1.45) circle(1.03);
 \fill[manifoldred] (4.33,.42) circle(3.7pt);
 \node[manifoldblue!85!black] at (4.33,1.45) {$e^1$};
 \draw[->,manifoldblue!75!black,thin]
  (4.62,1.74)--({4.33+1.03*cos(45)},{1.45+1.03*sin(45)});
 \node[below=7pt] at (4.33,.42) {$v$};
 \draw[->,manifoldgold!80!black,thick]
  (1.22,1.54)--(3.08,1.54)
  node[midway,above=5pt,font=\small] {dołącz $e^1$};
 % Dołączenie 2-komórki: brzeg dysku biegnie po okręgu.
 \fill[manifoldgreen!18] (8.68,1.45) circle(1.03);
 \draw[manifoldblue,line width=2.5pt]
  (8.68,1.45) circle(1.03);
 \fill[manifoldred] (8.68,.42) circle(3.7pt);
 \node[manifoldgreen!65!black] at (8.68,1.45) {$e^2$};
 \node[below=7pt] at (8.68,.42) {$v$};
 \draw[->,manifoldgold!80!black,thick]
  (5.58,1.54)--(7.35,1.54)
  node[midway,above=5pt,font=\small] {dołącz $e^2$};
\end{tikzpicture}
\caption{Przykłady kolejnych wymiarów CW: punkt
$X^0=\{v\}$, okrąg $X^1\cong S^1$ otrzymany
przez sklejenie obu końców $1$-komórki z $v$,
oraz dysk $X^2\cong D^2$ po dołączeniu
$2$-komórki wzdłuż całego okręgu.
Czerwony punkt, niebieski łuk i zielone wnętrze
oznaczają odpowiednio otwarte komórki
wymiarów $0,1,2$.}
\label{fig:cw-wymiary}
\end{figure}

W kompleksie skończonym warunki
topologiczne łatwo sprawdzić
bez osobnej teorii granic:
topologia słaba jest ilorazową
topologią sklejania skończenie
wielu dysków, a domknięcie
każdej komórki spotyka
tylko skończenie wiele innych.
Przy nieskończenie wielu komórkach
topologia słaba jest konieczna,
aby ciągłość dało się sprawdzać
na wszystkich domkniętych komórkach.

\begin{stwierdzenie}[Sympleksy dają komórki]
\label{prop:sympleksy-cw}
Skończony kompleks symplicjalny i
skończony $\Delta$-kompleks są
CW-kompleksami: ich otwarte komórki
to obrazy wnętrz sympleksów.
\end{stwierdzenie}
\begin{proof}
Indukujemy po wymiarze. Szkielet
zerowy składa się z wierzchołków.
Wnętrze każdego $n$-sympleksu
jest rozłączne z obrazami innych
wnętrz, a brzeg jest sumą ścian
zawartych w poprzednim szkielecie.
Stwierdzenie~\ref{prop:brzeg-sympleksu}
pozwala traktować sympleks
jako dysk $D^n$, którego
odwzorowanie na brzegu jest odwzorowaniem dołączającym.
Warunek o topologii jest w
Definicji~\ref{def:delta-kompleks}
podany wprost, a dla skończonej
realizacji symplicjalnej wynika
z topologii skończonej sumy zwartych sympleksów:
ciągła surjekcja z tej sumy na $|K|\subseteq\R^V$ jest domknięta,
gdyż obrazy zbiorów domkniętych są zwarte w przestrzeni Hausdorffa.
Jest więc odwzorowaniem ilorazowym, co daje właśnie wymaganą topologię.
Skończona liczba sympleksów
zapewnia skończoność domknięć komórek.
\end{proof}

\begin{przyklad}[Minimalne struktury komórkowe]
\label{prz:cw-sfera-torus}
\begin{enumerate}[label=(\arabic*),leftmargin=*]
\item Dla $n\geq1$ sfera $S^n$
      ma jedną $0$-komórkę i jedną
      $n$-komórkę:
      $D^n/S^{n-1}\cong S^n$.
      Jawny homeomorfizm opisujemy
      poniżej.
      Dla $S^0$ używamy dwóch
      $0$-komórek.
\item Torus ma jedną $0$-komórkę,
      dwie $1$-komórki $a,b$ oraz
      jedną $2$-komórkę. Obchodząc
      brzeg kwadratu z
      Rysunku~\ref{fig:delta-torus},
      otrzymujemy słowo
      $aba^{-1}b^{-1}$ w
      $S^1\vee S^1$.
      Dwa trójkąty realizacji
      $\Delta$-kompleksowej zostały
      zastąpione jednym dyskiem,
      którego odwzorowanie dołączające
      zapamiętuje sklejanie.
\end{enumerate}
\end{przyklad}

\begin{przyklad}[Sfera $S^2$ bez $1$-komórek]
\label{prz:cw-s2}
Weźmy $X^0=\{v\}$ i nie dołączajmy
żadnej $1$-komórki, więc $X^1=X^0$.
Niech odwzorowanie dołączające
$\varphi:\partial D^2=S^1\to\{v\}$
będzie stałe. Wtedy
\[
 X^2=X^1\cup_\varphi D^2
 \cong D^2/\partial D^2\cong S^2.
\]
Cały brzeg dysku staje się \emph{jednym}
punktem $v$; obraz jego wnętrza
jest jedną otwartą $2$-komórką
$e^2=S^2\setminus\{v\}\cong\R^2$.
Rysunek~\ref{fig:cw-s2} pokazuje to
sklejenie. W szczególności
$X^1$ nie jest okręgiem:
okrąg po lewej stronie rysunku
to brzeg \emph{dołączanego dysku},
a nie komórka przestrzeni wynikowej.
\end{przyklad}

\begin{figure}[!htbp]
\centering
\begin{tikzpicture}[>=Stealth,line cap=round,line join=round]
 % Dysk przed sklejeniem: czerwony brzeg zapada się do v.
 \fill[manifoldblue!10] (0,1.85) circle(1.33);
 \draw[manifoldred,line width=2.4pt] (0,1.85) circle(1.33);
 \node[font=\bfseries] at (0,3.58) {$D^2$};
 \node[manifoldblue!80!black] at (0,1.85) {$\operatorname{int}D^2$};
 \node[manifoldred!85!black,below=7pt]
  at (0,.52) {$\partial D^2=S^1$};
 % Sklejenie brzegu do jednego punktu.
 \draw[->,manifoldgold!85!black,line width=1.5pt]
  (1.73,1.85)--(4.8,1.85)
  node[midway,above=7pt] {$q$};
 \node[font=\small,manifoldgold!70!black]
  at (3.26,1.25) {sklejamy cały brzeg};
 % Sfera w stylistyce rysunku 1.4: jasny środek, tył i przód łuków.
 \shade[inner color=white,outer color=manifoldblue!43]
  (6.8,1.85) circle(1.43);
 \draw[manifoldblue!75!black,very thick]
  (6.8,1.85) circle(1.43);
 \draw[manifoldblue!50,densely dashed,thin]
  (8.23,1.85) arc[start angle=0,end angle=180,x radius=1.43,y radius=.34];
 \draw[manifoldblue!70!black,thin]
  (5.37,1.85) arc[start angle=180,end angle=360,x radius=1.43,y radius=.34];
 \draw[manifoldblue!38,densely dashed,thin]
  (6.8,3.28) .. controls (6.35,2.76) and (6.35,.94) .. (6.8,.42);
 \draw[manifoldblue!65!black,thin]
  (6.8,3.28) .. controls (7.25,2.76) and (7.25,.94) .. (6.8,.42);
 \fill[manifoldred] (6.8,.42) circle(3.8pt);
 \node[font=\bfseries] at (6.8,3.58) {$X^2\cong S^2$};
 \node[manifoldblue!85!black] at (6.8,2.44) {$e^2$};
 \node[manifoldred!85!black,below=7pt] at (6.8,.42) {$v$};
\end{tikzpicture}
\caption{Struktura CW sfery: $D^2$ dołączamy do
punktu $v$ odwzorowaniem stałym na całym brzegu;
$q:D^2\to D^2/\partial D^2$ jest mapą ilorazową.
W powstałej sferze czerwony punkt to
jedyne $e^0$, a cała pozostała powierzchnia
to jedna otwarta komórka $e^2$.
Nie występują komórki wymiaru $1$.}
\label{fig:cw-s2}
\end{figure}

\begin{uwaga}[Uwaga o ilorazie dysku]
W punkcie (1) nie potrzebujemy
identyfikować półkuli z ilorazem
za pomocą niejawnego rysunku.
Istnieje jawne ciągłe odwzorowanie
$D^n\to S^n\subseteq\R^{n+1}$,
\[
 x\longmapsto
 \bigl(\sin(\pi\|x\|)\,x/\|x\|,
       \cos(\pi\|x\|)\bigr)
 \quad(x\ne0),
\]
z wartością $(0,\ldots,0,1)$
w $x=0$. Cały brzeg $\|x\|=1$
przechodzi w biegun południowy,
a na wnętrzu odwzorowanie jest iniektywne
i pokrywa sferę bez tego bieguna.
Przechodzi więc na ciągłą bijekcję
$D^n/S^{n-1}\to S^n$, która jest
homeomorfizmem przez zwartość
dziedziny i własność Hausdorffa
sfery.
\end{uwaga}

\begin{twierdzenie}[Co robi dołączenie komórki z $\pi_1$?]
\label{thm:pi1-dolaczenie-komorki}
Niech $Y$ będzie łukowo spójna,
a $X=Y\cup_\varphi D^n$.
Jeśli $n=2$, po wyborze drogi od
punktu bazowego do $\varphi(S^1)$
grupa $\pi_1(X)$ jest ilorazem
$\pi_1(Y)$ przez domknięcie normalne
klasy pętli $\varphi$.
Jeśli $n\geq3$, włączenie
$Y\hookrightarrow X$ wyznacza
izomorfizm na $\pi_1$.
\end{twierdzenie}
\begin{proof}
W dysku wybierzmy promień $r=\|x\|$.
Otwarty w $X$ zbiór $U$ złożony
z $Y$ i pasa $r>1/3$ deformuje
się na $Y$: przesuwamy punkty
pasa promieniowo do $r=1$,
a punkty $Y$ zostawiamy. Na części dyskowej homotopia ma wzór
$ru\mapsto((1-t)r+t)u$, gdzie $u\in S^{n-1}$.
Wyjaśnijmy ciągłość przy sklejanym $Y$. Dla otwartego otoczenia
$O$ punktu $y\in Y$ zmniejszamy jego część dyskową do punktów $ru$,
dla których $\varphi(u)\in O\cap Y$ i cały odcinek
$\{\rho u:r\leq\rho\leq1\}$ leży w przeciwobrazie $O$.
Ta część jest otwarta: zawieranie zwartego odcinka w zbiorze
otwartym zachowuje się przy małej zmianie $r,u$.
Wraz z $O\cap Y$ tworzy otwarte otoczenie $y$ w ilorazie,
którego obraz przez homotopię pozostaje w $O$ dla wszystkich $t$.
Dowodzi to ciągłości przy $Y$; na wnętrzu dysku wynika ona z podanego wzoru.
Drugi otwarty zbiór $V$ jest
obrazem $r<2/3$; jest to
ściągalne wnętrze dysku.
Przecięcie $U\cap V$ to pierścień
$1/3<r<2/3$, homotopijnie
równoważny $S^{n-1}$.
Jest łukowo spójne dla $n\geq2$.
Otwartość obu zbiorów sprawdzamy przed sklejeniem:
przeciwobrazy $U$ to całe $Y$ i otwarty w $D^n$ pas $r>1/3$,
a przeciwobraz $V$ to $r<2/3$, bez punktów $Y$.

Twierdzenie~\ref{thm:van-kampen}
z $G_V=1$ mówi, że $\pi_1(X)$
jest ilorazem $\pi_1(U)\cong\pi_1(Y)$
przez domknięcie normalne obrazu
$\pi_1(U\cap V)$. Przy $n=2$
generator $\pi_1(S^1)\cong\Z$
po przejściu przez retrakcję
$U\to Y$ staje się pętlą
$\varphi$ (z uwzględnieniem
drogi do punktu bazowego).
Przy $n\geq3$ mamy
$\pi_1(S^{n-1})=0$ z
Przykładu~\ref{prz:pi1-sfera},
więc nie pojawia się żadna
relacja ani nowy generator.
\end{proof}

\begin{przyklad}[Torus i powierzchnia rzutowa]
\label{prz:pi1-torus-rp2}
Dla torusa część jednowymiarowa
jest $S^1\vee S^1$, z grupą
$F(a,b)$ z Przykładu~\ref{prz:pi1-dwa-okregi}.
Dołączona $2$-komórka wymusza
$aba^{-1}b^{-1}=1$, czyli
\[
 \pi_1(T^2)\cong
 \langle a,b\mid ab=ba\rangle
 \cong\Z^2.
\]
Drugim przykładem jest
$\R P^2$: dysk będący górną
półkulą $S^2$ ma na brzegu
antypody utożsamione, więc
$\R P^2$ ma $0$-komórkę,
      $1$-komórkę i $2$-komórkę
      dołączoną do $S^1\cong\R P^1$
      odwzorowaniem odpowiadającym
      $z\mapsto z^2$ (liczba owinięć wynosi $2$).
      Istotnie, $S^1/(z\sim-z)\to S^1$, $[z]\mapsto z^2$,
      jest homeomorfizmem jako ciągła bijekcja ze zwartej
      przestrzeni do przestrzeni Hausdorffa; odwzorowanie dołączające
      przechodzi przez ten iloraz. Pełny obieg brzegu daje zatem dwa obiegi
      $1$-szkieletu, co wyjaśnia relację $a^2=1$.
      Zatem
\[
 \pi_1(\R P^2)
 \cong\langle a\mid a^2=1\rangle
 \cong\Z/2\Z.
\]
W obu przykładach relacja grupowa pochodzi bezpośrednio
z przebiegu brzegu dołączanej $2$-komórki.
\end{przyklad}

\begin{uwaga}[Granica obecnego rozdziału]
Sympleksy i komórki dają
konkretne przestrzenie, a ich
odwzorowania dołączające kodują sposób
sklejenia. W następnym rozdziale
zorientowanym sympleksom przypiszemy
generatory łańcuchów, a ich ścianom
operator brzegu $\partial$.
Sprawdzimy $\partial^2=0$,
porównamy homologię symplicjalną,
singularną i komórkową oraz
wykorzystamy długie ciągi
z Rozdziału~\ref{ch:algebra-homologiczna}.
Wyższe grupy homotopii i twierdzenie
Hurewicza wymagają już tak
zbudowanej homologii przestrzeni.
\end{uwaga}

\section*{Dalsza lektura}
\addcontentsline{toc}{section}{Dalsza lektura}
Pełne rozwinięcie teorii homotopii,
nakryć, kompleksów CW i
$\Delta$-kompleksów podaje
Allen Hatcher, \emph{Algebraic Topology},
rozdziały 0 i 1 oraz początek rozdziału 2,
\url{https://pi.math.cornell.edu/~hatcher/AT/AT.pdf}.
W tym rozdziale wyprowadziliśmy
podstawowe własności, które
wykorzystamy do budowania homologii.

\clearpage
